\section{Experiments and Results}
\label{sec:experiments}

This section reports five experiments. The first three concern how student conflicts are measured
and whether the solver's objective can be tuned to reduce them; the fourth shows why several
apparently excellent results are not results at all; the fifth is the outcome of the phase.

\subsection{Solver checker versus KPI evaluator}

The first measurement to reconcile was a pair of figures that appeared to contradict each other on
the same term.

\begin{table}[H]
\centering
\caption{Two conflict rates for the Spr6 instance, and why they differ.}
\label{tab:reconcile}
\footnotesize
\begin{tabular}{@{}L{3.2cm}L{5.6cm}L{5.6cm}@{}}
\toprule
 & \textbf{Solver checker (Java)} & \textbf{KPI evaluator (CSV mode)} \\
\midrule
Value          & 526 of 990 conflicted, $C = 0.5313$ & 587 of 1{,}290 conflicted (45.5\,\%) \\
Instance       & \texttt{rl/instances/spr2026.xml} & \texttt{COEPSpr62026\_v2.csv} \\
Denominator    & students the solver loaded (990)  & all students in the enrolment file (1{,}290) \\
Membership     & the student block inside the solver file & re-derived by joining the CSV to file 16 \\
\bottomrule
\end{tabular}
\end{table}

The two are not inconsistent; they measure different populations of a different artefact. The
solver export carried enrolments for 990 of the 1{,}290 simulated students, and the CSV scored a
different exported solution. As Section~\ref{sec:evaluation} established, the two implementations
agree exactly when pointed at the same file. The lesson taken forward was procedural: every
conflict figure in this report is stated with its instance and its denominator, and no two figures
are compared unless both match.

\subsection{Tuning the student-conflict weight in UniTime}

CPSolver prices each student clash through \texttt{StudentConflictWeight}, and
\texttt{HardStudentConflictWeight} adds a further charge when the student has no alternative
section to move to. Raising the first should buy fewer conflicts at the cost of other soft
criteria. Two runs were made in the UniTime interface and exported as CSV.

\begin{table}[H]
\centering
\caption{Conflict weight raised in the UniTime solver, scored in CSV mode against the then-current
dataset (1{,}290 students, 262 classes).}
\label{tab:weightcsv}
\footnotesize
\begin{tabular}{@{}lrrrr@{}}
\toprule
\textbf{Weight} & \textbf{Students conflicted} & \textbf{Pairs} & \textbf{Hard feasibility} & \textbf{Preference score} \\
\midrule
5  & 676 of 1{,}290 & 1{,}076 & 90.97 & 25.00 \\
20 & 894 of 1{,}290 & 2{,}126 & 88.96 & 24.19 \\
\bottomrule
\end{tabular}
\end{table}

Read naively this says that raising the weight makes the timetable worse, which would be a strange
property for an optimiser. The correct reading is that this is not a controlled experiment. The
CSV join re-derives student membership from an enrolment file that the solver did not necessarily
use, the two runs started from different states, and the instance itself was changing during this
period. The experiment's real value was in motivating the controlled version that follows. The
figure of 676 at weight 5 supersedes an earlier figure of 614 quoted in the project's presentation
material, which no committed artefact reproduces.

\subsection{A controlled weight sweep}

To tune the weight properly the instance must be held fixed and the only variable must be the
weight. This was done headlessly, outside the UniTime server, using the project's own solver
runner. The instance was \texttt{rl/instances/spr2026.xml} (262 classes, 990 students, loaded with
an initial conflict rate of $C = 0.5313$); each configuration was run for 900 seconds of
wall-clock time with three random seeds; the reference arm B0 used the stock weight of 0.2.

\begin{figure}[H]
\centering
\begin{tikzpicture}
\begin{axis}[
  width=0.78\textwidth, height=6.2cm,
  xmode=log, log basis x=10,
  xlabel={\small \texttt{Comparator.StudentConflictWeight}},
  ylabel={\small student violation rate $C$},
  xtick={0.2,0.5,1,2,5,10}, xticklabels={0.2,0.5,1,2,5,10},
  ymin=-0.004, ymax=0.085,
  grid=major, grid style={gray!22},
  legend style={font=\scriptsize,at={(0.98,0.97)},anchor=north east,draw=gray!40},
  tick label style={font=\scriptsize}, label style={font=\small},
]
\addplot[only marks,mark=o,mark size=1.7pt,gray!65] table[x=weight,y=C] {data/sweep_points.dat};
\addlegendentry{individual seeds}
\addplot[coepblue,very thick,mark=*,mark size=2.2pt] table[x=weight,y=C] {data/sweep_median.dat};
\addlegendentry{median of 3 seeds}
\addplot[red!70!black,dashed,very thick,domain=0.18:11,samples=2] {0.04141};
\addlegendentry{B0 reference ($C=0.0414$)}
\end{axis}
\end{tikzpicture}
\caption{Student violation rate against the student-conflict weight, on a fixed instance with a
900-second budget and three seeds per weight. Every weight at or above 0.5 improves on the stock
B0 configuration, and seed-to-seed variance is comparable to the effect size.}
\label{fig:sweep}
\end{figure}

\begin{table}[H]
\centering
\caption{Weight sweep, median over three seeds (recorded in \texttt{rl/runs/sweep/summary.md}).}
\label{tab:sweep}
\footnotesize
\begin{tabular}{@{}lrrrr@{}}
\toprule
\textbf{Arm} & \textbf{Weight} & \textbf{$C$ (median)} & \textbf{Assigned} & \textbf{Soft penalty} \\
\midrule
B0 & 0.2  & 0.0414 & 1.00 & 60 \\
B1 & 0.2  & 0.0343 & 1.00 & 62 \\
B1 & 0.5  & 0.0192 & 1.00 & 66 \\
B1 & 1    & 0.0263 & 1.00 & 67 \\
B1 & 2    & 0.0030 & 1.00 & 65 \\
B1 & 5    & 0.0000 & 1.00 & 69 \\
B1 & 10   & 0.0040 & 1.00 & 64 \\
\bottomrule
\end{tabular}
\end{table}

Two observations follow. First, the solver starts from $C = 0.5313$ and every configuration drives
it below 0.06 within the budget while keeping all 262 classes assigned, so the overwhelming
majority of the reduction comes from the search itself rather than from the weight. Second, the
weight does matter at the margin: the best setting reduces the residual rate by roughly an order of
magnitude relative to the stock value.

One caveat must be attached to the weight-5 row. Recomputing the medians from the per-seed metric
files currently on disk reproduces every row of Table~\ref{tab:sweep} except that one, for which
the stored files give a median of 0.0263 and make weight 2 ($C = 0.0030$) the minimum instead. The
discrepancy is confined to a single seed whose run appears to have been repeated after the summary
was written. The choice between the two readings does not affect any conclusion: measured against
B0's 0.0414, the best weight achieves a relative reduction of either 1.000 or 0.927, and both
clear the 25\,\% threshold that had been fixed in advance. The honest statement is that weight
tuning alone reduces the residual conflict rate by well over 25\,\%, not that it reaches exactly
zero.

\subsection{Why several ``zero conflict'' results are not results}

Four artefacts in the repository report zero student conflicts. None of them represents a better
timetable, and Table~\ref{tab:zeros} gives the mechanism in each case.

\begin{table}[H]
\centering
\caption{Zero-conflict results and their actual causes.}
\label{tab:zeros}
\footnotesize
\begin{tabular}{@{}L{4.6cm}L{3.0cm}L{6.8cm}@{}}
\toprule
\textbf{Artefact} & \textbf{Reported} & \textbf{Mechanism} \\
\midrule
\texttt{spr2026\_mihir.kpi.json}, \texttt{spr2026\_10m\_w5.kpi.json}
  & 0 of 990
  & Student switching and initial sectioning moved students between parallel laboratory batches
    until the overlaps vanished. The placements did not change; the cohort structure did. \\
\texttt{rl/runs/bakeoff-spr2026\_10m\_w5} (3 arms $\times$ 10 seeds, 900\,s)
  & 0 for all arms
  & The instance was \emph{already} conflict-free when loaded: the initial violation rate recorded
    in every metrics file is 0.0. A 900-second run that begins and ends at zero measures nothing
    about conflict reduction. \\
\texttt{spr2026\_10\_LE\_90m\_w5.kpi.json}
  & 0 of 0
  & No students were loaded at all. The rate is zero over an empty denominator. \\
\texttt{kpi\_report\_MacMini\_v1.json}
  & 0 of 2{,}130
  & A different and earlier instance with only 254 classes, and two instructor conflicts. \\
\bottomrule
\end{tabular}
\end{table}

The first row is the mechanism that mattered. UniTime's final sectioning pass is permitted, by
default, to move a student from one laboratory batch to another in order to resolve a clash. For a
university that assigns batches administratively this is not a solution but a corruption of the
input, and it makes the reported metric unfalsifiable: any timetable can be made conflict-free if
the students are allowed to move. The remedy is the three settings of Section~\ref{sec:setup}, and
applying them is what produced the final experiment.

\subsection{The final solution under pinned enrolments}

With \emph{Enrolled Student Course Demands} selected, \texttt{General.SwitchStudents} set to
\texttt{false} and \texttt{Global.LoadStudentEnrlsFromSolution} set to \texttt{true}, the solver
can no longer move students, and the conflict count becomes a property of the timetable alone. The
exported solution \texttt{spr2026\_11\_all\_flags.csv} was scored against the current dataset:

\begin{lstlisting}[style=plainlst,language=bash]
python scripts/eval_timetable_kpis.py \
    --csv unitime-out/spr2026_11_all_flags.csv \
    --data-dir unitime-out/0gen_all
\end{lstlisting}

\begin{figure}[H]
\centering
\begin{tikzpicture}
\begin{axis}[
  width=0.95\textwidth, height=6.4cm,
  ybar, bar width=15pt,
  ymin=85, ymax=101.5,
  ylabel={\small pass rate (\%)},
  symbolic x coords={Coverage,Time pattern,Date pattern,Room capacity,Lunch/weekend,Instr. lunch,Student lunch,Instr. clash,Room clash,Student conflicts},
  xtick=data,
  x tick label style={rotate=38,anchor=east,font=\scriptsize},
  tick label style={font=\scriptsize}, label style={font=\small},
  nodes near coords, every node near coord/.append style={font=\tiny,rotate=90,anchor=west},
  grid=major, grid style={gray!20},
  enlarge x limits=0.055,
]
\addplot[fill=coepblue!55,draw=coepblue] coordinates {
 (Coverage,99.66) (Time pattern,100) (Date pattern,100) (Room capacity,100)
 (Lunch/weekend,100) (Instr. lunch,97.92) (Student lunch,94.42)
 (Instr. clash,100) (Room clash,100) (Student conflicts,96.46)};
\end{axis}
\end{tikzpicture}
\caption{Hard-check pass rates of the final solution (term Spr11 solution scored against the
Spr12 dataset). The vertical axis starts at 85\,\% to make the three non-perfect checks legible.}
\label{fig:hardchecks}
\end{figure}

\begin{table}[H]
\centering
\caption{Full evaluation of the final solution: 291 classes, 1{,}130 students, 42 rooms,
115 staff.}
\label{tab:final}
\footnotesize
\begin{tabular}{@{}L{6.6cm}L{4.0cm}r@{}}
\toprule
\textbf{Metric} & \textbf{Value} & \textbf{Rate} \\
\midrule
Classes in offerings / in solution    & 291 / 291        & --- \\
Coverage                              & 290 of 291 assigned & 99.66\,\% \\
Hard feasibility score                & \textbf{98.85} / 100 & --- \\
Hard checks / failures                & 4{,}341 / 105    & --- \\
Time-pattern fidelity                 & 290 of 290       & 100\,\% \\
Date pattern                          & 290 of 290       & 100\,\% \\
Room capacity                         & 290 of 290       & 100\,\% \\
Lunch/weekend slot compliance         & 290 of 290       & 100\,\% \\
Instructor double-booking             & 0 conflicts      & 100\,\% \\
Room double-booking                   & 0 conflicts      & 100\,\% \\
\textbf{Student conflicts}            & \textbf{40 of 1{,}130 students, 40 pairs} & \textbf{96.46\,\%} \\
Student lunch breaks                  & 63 violations    & 94.42\,\% \\
Instructor lunch breaks               & 1 violation      & 97.92\,\% \\
Travel-time violations                & 0                & --- \\
Back-to-back density                  & 31.96\,\%        & --- \\
Unavailable/online room usage         & 17 of 291        & 5.84\,\% \\
Preference score                      & 6.29 / 100       & --- \\
\quad room preference                 & 35 of 278        & 12.59\,\% \\
\quad instructor identity             & 0 of 289 (artefact) & 0\,\% \\
\bottomrule
\end{tabular}
\end{table}

\paragraph{Interpretation.} The timetable is hard-feasible in every respect that the institution
treats as non-negotiable: no room is double-booked, no instructor is double-booked, every class
respects its meeting pattern and its room capacity, and 290 of 291 classes are placed. The
remaining student conflicts are a small and highly localised set. All 40 are first-year students,
and all 40 arise from the same pair of courses: a first-year programming laboratory overlapping the
first-year reserved lecture block, in four specific lab/lecture combinations---laboratory 1 against
reserved lecture 3 on Tuesday, laboratory 2 against reserved lecture 3 on Friday, laboratory 7
against reserved lecture 2 on Wednesday, and laboratory 9 against reserved lecture 1 or 2 on
Friday. Because the conflicts are concentrated in four placements rather than spread across the
cohort, they are addressable either by re-timing those four laboratories or by adding explicit
different-time distribution constraints between the laboratory and the reserved block.

Two results should not be over-read. The room preference rate of 12.59\,\% is genuinely poor and is
the clearest remaining optimisation target; it reflects that room preferences were generated
permissively from Taasika associations and that the solver traded them away cheaply. The
instructor identity rate of 0\,\% is not a result at all, but the name-format artefact described in
Section~\ref{sec:evaluation}.

\subsection{Cross-run comparison}

Table~\ref{tab:crossrun} collects every evaluated artefact so that no figure in this report
circulates without its context. The progression from the top group to the bottom row is the
substance of the phase: the apparent zeros disappear as the measurement is tightened, and what
remains is a small, real, diagnosable number.

\begin{table}[H]
\centering
\caption{All evaluated runs. CSV mode joins a timetable to an enrolment file; XML mode scores the
student block inside a solver file.}
\label{tab:crossrun}
\footnotesize
\begin{tabular}{@{}L{4.9cm}llrrr@{}}
\toprule
\textbf{Artefact} & \textbf{Mode} & \textbf{Instance} & \textbf{Classes} & \textbf{Conflicted} & \textbf{Hard feas.} \\
\midrule
\texttt{spr2026.kpi.json}               & XML & Spr6  & 262 & 526 / 990   & 92.15 \\
\texttt{kpi\_report.json}               & CSV & Spr6  & 260 & 587 / 1{,}290 & 94.98 \\
\texttt{kpi\_report\_2785.json}         & CSV & Spr6  & 262 & 507 / 1{,}290 & 92.66 \\
\texttt{kpi\_weight5\_report.json}      & CSV & Spr6  & 262 & 676 / 1{,}290 & 90.97 \\
\texttt{kpi\_weight20\_report.json}     & CSV & Spr6  & 262 & 894 / 1{,}290 & 88.96 \\
\texttt{solution.kpi.json} (sweep, $w{=}5$) & XML & Spr6 & 262 & 26 / 990 & 97.42 \\
\texttt{spr2026\_mihir.kpi.json}        & XML & Spr6  & 262 & 0 / 990 \textsuperscript{a} & 97.67 \\
\texttt{spr2026\_10m\_w5.kpi.json}      & XML & Spr6  & 262 & 0 / 990 \textsuperscript{a} & 98.26 \\
\texttt{spr2026\_10\_LE\_90m\_w5}       & XML & Spr10 & 290 & 0 / 0 \textsuperscript{b}   & 97.55 \\
\texttt{spr2026\_10\_LE\_SS\_Enrolled}  & XML & Spr10 & 290 & 32 / 1{,}130 & 97.77 \\
\textbf{\texttt{spr2026\_11\_all\_flags.csv}} & CSV & Spr11/12 & 291 & \textbf{40 / 1{,}130} & \textbf{98.85} \\
\bottomrule
\multicolumn{6}{@{}l@{}}{\footnotesize \textsuperscript{a} students re-sectioned;
  \textsuperscript{b} no students loaded. Neither is a conflict-reduction result.}
\end{tabular}
\end{table}
